\IfFileExists{stacks-project.cls}{%
\documentclass{stacks-project}
}{%
\documentclass{amsart}
}

% For dealing with references we use the comment environment
\usepackage{verbatim}
\newenvironment{reference}{\comment}{\endcomment}
%\newenvironment{reference}{}{}
\newenvironment{slogan}{\comment}{\endcomment}
\newenvironment{history}{\comment}{\endcomment}

% For commutative diagrams we use Xy-pic
\usepackage[all]{xy}

% We use 2cell for 2-commutative diagrams.
\xyoption{2cell}
\UseAllTwocells

% We use multicol for the list of chapters between chapters
\usepackage{multicol}

% This is generally recommended for better output
\usepackage{lmodern}
\usepackage[T1]{fontenc}

% For cross-file-references
\usepackage{xr-hyper}

% Package for hypertext links:
\usepackage{hyperref}

% For any local file, say "hello.tex" you want to link to please
% use \externaldocument[hello-]{hello}
\externaldocument[introduction-]{introduction}
\externaldocument[conventions-]{conventions}
\externaldocument[sets-]{sets}
\externaldocument[categories-]{categories}
\externaldocument[topology-]{topology}
\externaldocument[sheaves-]{sheaves}
\externaldocument[sites-]{sites}
\externaldocument[stacks-]{stacks}
\externaldocument[fields-]{fields}
\externaldocument[algebra-]{algebra}
\externaldocument[brauer-]{brauer}
\externaldocument[homology-]{homology}
\externaldocument[derived-]{derived}
\externaldocument[simplicial-]{simplicial}
\externaldocument[more-algebra-]{more-algebra}
\externaldocument[smoothing-]{smoothing}
\externaldocument[modules-]{modules}
\externaldocument[sites-modules-]{sites-modules}
\externaldocument[injectives-]{injectives}
\externaldocument[cohomology-]{cohomology}
\externaldocument[sites-cohomology-]{sites-cohomology}
\externaldocument[dga-]{dga}
\externaldocument[dpa-]{dpa}
\externaldocument[sdga-]{sdga}
\externaldocument[hypercovering-]{hypercovering}
\externaldocument[schemes-]{schemes}
\externaldocument[constructions-]{constructions}
\externaldocument[properties-]{properties}
\externaldocument[morphisms-]{morphisms}
\externaldocument[coherent-]{coherent}
\externaldocument[divisors-]{divisors}
\externaldocument[limits-]{limits}
\externaldocument[varieties-]{varieties}
\externaldocument[topologies-]{topologies}
\externaldocument[descent-]{descent}
\externaldocument[perfect-]{perfect}
\externaldocument[more-morphisms-]{more-morphisms}
\externaldocument[flat-]{flat}
\externaldocument[groupoids-]{groupoids}
\externaldocument[more-groupoids-]{more-groupoids}
\externaldocument[etale-]{etale}
\externaldocument[chow-]{chow}
\externaldocument[intersection-]{intersection}
\externaldocument[pic-]{pic}
\externaldocument[weil-]{weil}
\externaldocument[adequate-]{adequate}
\externaldocument[dualizing-]{dualizing}
\externaldocument[duality-]{duality}
\externaldocument[discriminant-]{discriminant}
\externaldocument[derham-]{derham}
\externaldocument[local-cohomology-]{local-cohomology}
\externaldocument[algebraization-]{algebraization}
\externaldocument[curves-]{curves}
\externaldocument[resolve-]{resolve}
\externaldocument[models-]{models}
\externaldocument[functors-]{functors}
\externaldocument[equiv-]{equiv}
\externaldocument[pione-]{pione}
\externaldocument[etale-cohomology-]{etale-cohomology}
\externaldocument[proetale-]{proetale}
\externaldocument[relative-cycles-]{relative-cycles}
\externaldocument[more-etale-]{more-etale}
\externaldocument[trace-]{trace}
\externaldocument[crystalline-]{crystalline}
\externaldocument[spaces-]{spaces}
\externaldocument[spaces-properties-]{spaces-properties}
\externaldocument[spaces-morphisms-]{spaces-morphisms}
\externaldocument[decent-spaces-]{decent-spaces}
\externaldocument[spaces-cohomology-]{spaces-cohomology}
\externaldocument[spaces-limits-]{spaces-limits}
\externaldocument[spaces-divisors-]{spaces-divisors}
\externaldocument[spaces-over-fields-]{spaces-over-fields}
\externaldocument[spaces-topologies-]{spaces-topologies}
\externaldocument[spaces-descent-]{spaces-descent}
\externaldocument[spaces-perfect-]{spaces-perfect}
\externaldocument[spaces-more-morphisms-]{spaces-more-morphisms}
\externaldocument[spaces-flat-]{spaces-flat}
\externaldocument[spaces-groupoids-]{spaces-groupoids}
\externaldocument[spaces-more-groupoids-]{spaces-more-groupoids}
\externaldocument[bootstrap-]{bootstrap}
\externaldocument[spaces-pushouts-]{spaces-pushouts}
\externaldocument[spaces-chow-]{spaces-chow}
\externaldocument[groupoids-quotients-]{groupoids-quotients}
\externaldocument[spaces-more-cohomology-]{spaces-more-cohomology}
\externaldocument[spaces-simplicial-]{spaces-simplicial}
\externaldocument[spaces-duality-]{spaces-duality}
\externaldocument[formal-spaces-]{formal-spaces}
\externaldocument[restricted-]{restricted}
\externaldocument[spaces-resolve-]{spaces-resolve}
\externaldocument[formal-defos-]{formal-defos}
\externaldocument[defos-]{defos}
\externaldocument[cotangent-]{cotangent}
\externaldocument[examples-defos-]{examples-defos}
\externaldocument[algebraic-]{algebraic}
\externaldocument[examples-stacks-]{examples-stacks}
\externaldocument[stacks-sheaves-]{stacks-sheaves}
\externaldocument[criteria-]{criteria}
\externaldocument[artin-]{artin}
\externaldocument[quot-]{quot}
\externaldocument[stacks-properties-]{stacks-properties}
\externaldocument[stacks-morphisms-]{stacks-morphisms}
\externaldocument[stacks-limits-]{stacks-limits}
\externaldocument[stacks-cohomology-]{stacks-cohomology}
\externaldocument[stacks-perfect-]{stacks-perfect}
\externaldocument[stacks-introduction-]{stacks-introduction}
\externaldocument[stacks-more-morphisms-]{stacks-more-morphisms}
\externaldocument[stacks-geometry-]{stacks-geometry}
\externaldocument[moduli-]{moduli}
\externaldocument[moduli-curves-]{moduli-curves}
\externaldocument[examples-]{examples}
\externaldocument[exercises-]{exercises}
\externaldocument[guide-]{guide}
\externaldocument[desirables-]{desirables}
\externaldocument[coding-]{coding}
\externaldocument[obsolete-]{obsolete}
\externaldocument[fdl-]{fdl}
\externaldocument[index-]{index}

% Theorem environments.
%
\theoremstyle{plain}
\newtheorem{theorem}[subsection]{Theorem}
\newtheorem{proposition}[subsection]{Proposition}
\newtheorem{lemma}[subsection]{Lemma}

\theoremstyle{definition}
\newtheorem{definition}[subsection]{Definition}
\newtheorem{example}[subsection]{Example}
\newtheorem{exercise}[subsection]{Exercise}
\newtheorem{situation}[subsection]{Situation}

\theoremstyle{remark}
\newtheorem{remark}[subsection]{Remark}
\newtheorem{remarks}[subsection]{Remarks}

\numberwithin{equation}{subsection}

% Macros
%
\def\lim{\mathop{\mathrm{lim}}\nolimits}
\def\colim{\mathop{\mathrm{colim}}\nolimits}
\def\Spec{\mathop{\mathrm{Spec}}}
\def\Hom{\mathop{\mathrm{Hom}}\nolimits}
\def\Ext{\mathop{\mathrm{Ext}}\nolimits}
\def\SheafHom{\mathop{\mathcal{H}\!\mathit{om}}\nolimits}
\def\SheafExt{\mathop{\mathcal{E}\!\mathit{xt}}\nolimits}
\def\Sch{\mathit{Sch}}
\def\Mor{\mathop{\mathrm{Mor}}\nolimits}
\def\Ob{\mathop{\mathrm{Ob}}\nolimits}
\def\Sh{\mathop{\mathit{Sh}}\nolimits}
\def\NL{\mathop{N\!L}\nolimits}
\def\CH{\mathop{\mathrm{CH}}\nolimits}
\def\proetale{{pro\text{-}\acute{e}tale}}
\def\etale{{\acute{e}tale}}
\def\QCoh{\mathit{QCoh}}
\def\Ker{\mathop{\mathrm{Ker}}}
\def\Im{\mathop{\mathrm{Im}}}
\def\Coker{\mathop{\mathrm{Coker}}}
\def\Coim{\mathop{\mathrm{Coim}}}

% Boxtimes
%
\DeclareMathSymbol{\boxtimes}{\mathbin}{AMSa}{"02}

%
% Macros for moduli stacks/spaces
%
\def\QCohstack{\mathcal{QC}\!\mathit{oh}}
\def\Cohstack{\mathcal{C}\!\mathit{oh}}
\def\Spacesstack{\mathcal{S}\!\mathit{paces}}
\def\Quotfunctor{\mathrm{Quot}}
\def\Hilbfunctor{\mathrm{Hilb}}
\def\Curvesstack{\mathcal{C}\!\mathit{urves}}
\def\Polarizedstack{\mathcal{P}\!\mathit{olarized}}
\def\Complexesstack{\mathcal{C}\!\mathit{omplexes}}
% \Pic is the operator that assigns to X its picard group, usage \Pic(X)
% \Picardstack_{X/B} denotes the Picard stack of X over B
% \Picardfunctor_{X/B} denotes the Picard functor of X over B
\def\Pic{\mathop{\mathrm{Pic}}\nolimits}
\def\Picardstack{\mathcal{P}\!\mathit{ic}}
\def\Picardfunctor{\mathrm{Pic}}
\def\Deformationcategory{\mathcal{D}\!\mathit{ef}}
\usepackage{cleveref}
\crefname{originalreference}{original reference}{original references}
\Crefname{originalreference}{Original reference}{Original references}
\makeatletter
\newlabel{algebra-definition-separable-field-extension}{{030O}{source}{}{sourcepointer-1}{}}
\newlabel{algebra-definition-separable-field-extension@cref}{{[originalreference][0][]030O}{[1][0][]source}{}{}{}}
\makeatother
\crefname{originalreference}{original reference}{original references}
\Crefname{originalreference}{Original reference}{Original references}
\makeatletter
\newlabel{algebra-definition-smooth}{{00T2}{source}{}{sourcepointer-2}{}}
\newlabel{algebra-definition-smooth@cref}{{[originalreference][0][]00T2}{[1][0][]source}{}{}{}}
\makeatother
\crefname{originalreference}{original reference}{original references}
\Crefname{originalreference}{Original reference}{Original references}
\makeatletter
\newlabel{algebra-lemma-NL-surjection}{{07BP}{source}{}{sourcepointer-3}{}}
\newlabel{algebra-lemma-NL-surjection@cref}{{[originalreference][0][]07BP}{[1][0][]source}{}{}{}}
\makeatother
\crefname{originalreference}{original reference}{original references}
\Crefname{originalreference}{Original reference}{Original references}
\makeatletter
\newlabel{algebra-lemma-base-change-smooth}{{00T4}{source}{}{sourcepointer-4}{}}
\newlabel{algebra-lemma-base-change-smooth@cref}{{[originalreference][0][]00T4}{[1][0][]source}{}{}{}}
\makeatother
\crefname{originalreference}{original reference}{original references}
\Crefname{originalreference}{Original reference}{Original references}
\makeatletter
\newlabel{algebra-lemma-characterize-smooth-over-field}{{00TT}{source}{}{sourcepointer-5}{}}
\newlabel{algebra-lemma-characterize-smooth-over-field@cref}{{[originalreference][0][]00TT}{[1][0][]source}{}{}{}}
\makeatother
\crefname{originalreference}{original reference}{original references}
\Crefname{originalreference}{Original reference}{Original references}
\makeatletter
\newlabel{algebra-lemma-computation-differential}{{00TU}{source}{}{sourcepointer-6}{}}
\newlabel{algebra-lemma-computation-differential@cref}{{[originalreference][0][]00TU}{[1][0][]source}{}{}{}}
\makeatother
\crefname{originalreference}{original reference}{original references}
\Crefname{originalreference}{Original reference}{Original references}
\makeatletter
\newlabel{algebra-lemma-differentials-base-change}{{00RV}{source}{}{sourcepointer-7}{}}
\newlabel{algebra-lemma-differentials-base-change@cref}{{[originalreference][0][]00RV}{[1][0][]source}{}{}{}}
\makeatother
\crefname{originalreference}{original reference}{original references}
\Crefname{originalreference}{Original reference}{Original references}
\makeatletter
\newlabel{algebra-lemma-exact-sequence-NL}{{00S2}{source}{}{sourcepointer-8}{}}
\newlabel{algebra-lemma-exact-sequence-NL@cref}{{[originalreference][0][]00S2}{[1][0][]source}{}{}{}}
\makeatother
\crefname{originalreference}{original reference}{original references}
\Crefname{originalreference}{Original reference}{Original references}
\makeatletter
\newlabel{algebra-lemma-regular-ring-CM}{{00NQ}{source}{}{sourcepointer-9}{}}
\newlabel{algebra-lemma-regular-ring-CM@cref}{{[originalreference][0][]00NQ}{[1][0][]source}{}{}{}}
\makeatother
\crefname{originalreference}{original reference}{original references}
\Crefname{originalreference}{Original reference}{Original references}
\makeatletter
\newlabel{algebra-lemma-regular-sequence-in-neighbourhood}{{061L}{source}{}{sourcepointer-10}{}}
\newlabel{algebra-lemma-regular-sequence-in-neighbourhood@cref}{{[originalreference][0][]061L}{[1][0][]source}{}{}{}}
\makeatother
\crefname{originalreference}{original reference}{original references}
\Crefname{originalreference}{Original reference}{Original references}
\makeatletter
\newlabel{algebra-lemma-smooth-at-point}{{07BU}{source}{}{sourcepointer-11}{}}
\newlabel{algebra-lemma-smooth-at-point@cref}{{[originalreference][0][]07BU}{[1][0][]source}{}{}{}}
\makeatother
\crefname{originallemma}{original lemma}{original lemmas}
\Crefname{originallemma}{Original lemma}{Original lemmas}
\makeatletter
\newlabel{lemma-neron-functorial}{{0BJ3}{source}{}{sourcepointer-12}{}}
\newlabel{lemma-neron-functorial@cref}{{[originallemma][0][]0BJ3}{[1][0][]source}{}{}{}}
\makeatother
\crefname{originalreference}{original reference}{original references}
\Crefname{originalreference}{Original reference}{Original references}
\makeatletter
\newlabel{more-algebra-lemma-blowup-regular-sequence}{{0BIQ}{source}{}{sourcepointer-13}{}}
\newlabel{more-algebra-lemma-blowup-regular-sequence@cref}{{[originalreference][0][]0BIQ}{[1][0][]source}{}{}{}}
\makeatother
\crefname{originalreference}{original reference}{original references}
\Crefname{originalreference}{Original reference}{Original references}
\makeatletter
\newlabel{situation-neron}{{0BJ2}{source}{}{sourcepointer-14}{}}
\newlabel{situation-neron@cref}{{[originalreference][0][]0BJ2}{[1][0][]source}{}{}{}}
\makeatother
\title{Neron's blowup of a smooth-at-p algebra is smooth at p' with an explicit cotangent sequence}
\author{The Stacks Project Authors}
\date{stacks-project:smoothing.tex:local-snapshot-20261012; GNU FDL 1.2 or later}
\begin{document}
\maketitle
\noindent\textit{Editorial note.} Input lemma-neron-functorial is same-chapter NOT carded (dropped in judgment as a pure pointer to Algebra blowup-base-change) and is declared external.\par\medskip
\noindent\textit{Allowed standard prerequisites (editorial).}\begin{itemize}
\item Functoriality of Neron's blowup. lemma-neron-functorial (same chapter, NOT carded, declared external). Permits replacing A by localizations off p throughout the proof.
\item Base change of smooth ring maps. Algebra, lemma base-change-smooth (cross-chapter, named). A/pi A is smooth over kappa at p.
\item Blowup along a regular sequence. More on Algebra, lemma blowup-regular-sequence (cross-chapter, named). Gives the presentation A' = A[y\_1,...,y\_c]/(pi y\_i - g\_i).
\item Naive cotangent exact sequence for a surjection. Algebra, lemma exact-sequence-NL and lemma NL-surjection (cross-chapter, named). Computes H\_1(NL\_\{A'/R\}) and the conormal map of the blowup presentation.
\item Independence of differentials of parameters. Algebra, lemma computation-differential (cross-chapter, named). Gives the kappa(p)-linear independence of dg\_i from separability of the residue extension.
\end{itemize}\medskip
\ifdefined\hypertarget\hypertarget{sourcepointer-1}{}\fi
\noindent\textit{Source reference.} Original reference 030O: focused-seeds/geometry/S9-smo-neron-blowup-smooth/ref-tags.txt, source lines 5--5. Navigation pointer only: the counted proof cites this statement by \textbackslash{}ref as a standing input; the deduction the counted proof has to make is carried out inside the counted span, and no part of that deduction is delegated to this pointer.\par\medskip
\ifdefined\hypertarget\hypertarget{sourcepointer-2}{}\fi
\noindent\textit{Source reference.} Original reference 00T2: focused-seeds/geometry/S9-smo-neron-blowup-smooth/ref-tags.txt, source lines 6--6. Navigation pointer only: the counted proof cites this statement by \textbackslash{}ref as a standing input; the deduction the counted proof has to make is carried out inside the counted span, and no part of that deduction is delegated to this pointer.\par\medskip
\ifdefined\hypertarget\hypertarget{sourcepointer-3}{}\fi
\noindent\textit{Source reference.} Original reference 07BP: focused-seeds/geometry/S9-smo-neron-blowup-smooth/ref-tags.txt, source lines 7--7. Navigation pointer only: the counted proof cites this statement by \textbackslash{}ref as a standing input; the deduction the counted proof has to make is carried out inside the counted span, and no part of that deduction is delegated to this pointer.\par\medskip
\ifdefined\hypertarget\hypertarget{sourcepointer-4}{}\fi
\noindent\textit{Source reference.} Original reference 00T4: focused-seeds/geometry/S9-smo-neron-blowup-smooth/ref-tags.txt, source lines 8--8. Navigation pointer only: the counted proof cites this statement by \textbackslash{}ref as a standing input; the deduction the counted proof has to make is carried out inside the counted span, and no part of that deduction is delegated to this pointer.\par\medskip
\ifdefined\hypertarget\hypertarget{sourcepointer-5}{}\fi
\noindent\textit{Source reference.} Original reference 00TT: focused-seeds/geometry/S9-smo-neron-blowup-smooth/ref-tags.txt, source lines 9--9. Navigation pointer only: the counted proof cites this statement by \textbackslash{}ref as a standing input; the deduction the counted proof has to make is carried out inside the counted span, and no part of that deduction is delegated to this pointer.\par\medskip
\ifdefined\hypertarget\hypertarget{sourcepointer-6}{}\fi
\noindent\textit{Source reference.} Original reference 00TU: focused-seeds/geometry/S9-smo-neron-blowup-smooth/ref-tags.txt, source lines 10--10. Navigation pointer only: the counted proof cites this statement by \textbackslash{}ref as a standing input; the deduction the counted proof has to make is carried out inside the counted span, and no part of that deduction is delegated to this pointer.\par\medskip
\ifdefined\hypertarget\hypertarget{sourcepointer-7}{}\fi
\noindent\textit{Source reference.} Original reference 00RV: focused-seeds/geometry/S9-smo-neron-blowup-smooth/ref-tags.txt, source lines 11--11. Navigation pointer only: the counted proof cites this statement by \textbackslash{}ref as a standing input; the deduction the counted proof has to make is carried out inside the counted span, and no part of that deduction is delegated to this pointer.\par\medskip
\ifdefined\hypertarget\hypertarget{sourcepointer-8}{}\fi
\noindent\textit{Source reference.} Original reference 00S2: focused-seeds/geometry/S9-smo-neron-blowup-smooth/ref-tags.txt, source lines 12--12. Navigation pointer only: the counted proof cites this statement by \textbackslash{}ref as a standing input; the deduction the counted proof has to make is carried out inside the counted span, and no part of that deduction is delegated to this pointer.\par\medskip
\ifdefined\hypertarget\hypertarget{sourcepointer-9}{}\fi
\noindent\textit{Source reference.} Original reference 00NQ: focused-seeds/geometry/S9-smo-neron-blowup-smooth/ref-tags.txt, source lines 13--13. Navigation pointer only: the counted proof cites this statement by \textbackslash{}ref as a standing input; the deduction the counted proof has to make is carried out inside the counted span, and no part of that deduction is delegated to this pointer.\par\medskip
\ifdefined\hypertarget\hypertarget{sourcepointer-10}{}\fi
\noindent\textit{Source reference.} Original reference 061L: focused-seeds/geometry/S9-smo-neron-blowup-smooth/ref-tags.txt, source lines 14--14. Navigation pointer only: the counted proof cites this statement by \textbackslash{}ref as a standing input; the deduction the counted proof has to make is carried out inside the counted span, and no part of that deduction is delegated to this pointer.\par\medskip
\ifdefined\hypertarget\hypertarget{sourcepointer-11}{}\fi
\noindent\textit{Source reference.} Original reference 07BU: focused-seeds/geometry/S9-smo-neron-blowup-smooth/ref-tags.txt, source lines 15--15. Navigation pointer only: the counted proof cites this statement by \textbackslash{}ref as a standing input; the deduction the counted proof has to make is carried out inside the counted span, and no part of that deduction is delegated to this pointer.\par\medskip
\ifdefined\hypertarget\hypertarget{sourcepointer-12}{}\fi
\noindent\textit{Source reference.} Original lemma 0BJ3: sources/stacks-smoothing/pinned-original-source.tex, source lines 1124--1124. Navigation pointer only: the counted proof cites this statement by \textbackslash{}ref as a standing input; the deduction the counted proof has to make is carried out inside the counted span, and no part of that deduction is delegated to this pointer.\par\medskip
\ifdefined\hypertarget\hypertarget{sourcepointer-13}{}\fi
\noindent\textit{Source reference.} Original reference 0BIQ: focused-seeds/geometry/S9-smo-neron-blowup-smooth/ref-tags.txt, source lines 16--16. Navigation pointer only: the counted proof cites this statement by \textbackslash{}ref as a standing input; the deduction the counted proof has to make is carried out inside the counted span, and no part of that deduction is delegated to this pointer.\par\medskip
\ifdefined\hypertarget\hypertarget{sourcepointer-14}{}\fi
\noindent\textit{Source reference.} Original reference 0BJ2: sources/stacks-smoothing/pinned-original-source.tex, source lines 1085--1085. Navigation pointer only: the counted proof cites this statement by \textbackslash{}ref as a standing input; the deduction the counted proof has to make is carried out inside the counted span, and no part of that deduction is delegated to this pointer.\par\medskip
\begin{lemma}
\label{lemma-neron-blowup-smooth}
In Situation \ref{situation-neron} assume that $R \to A$ is smooth
at $\mathfrak p$ and that $R/\pi R \subset \Lambda/\pi \Lambda$
is a separable field extension. Then $R \to A'$ is smooth at
$\mathfrak p'$ and there is a short exact sequence
$$
0 \to
\Omega_{A/R} \otimes_A A'_{\mathfrak p'} \to
\Omega_{A'/R, \mathfrak p'} \to
(A'/\pi A')_{\mathfrak p'}^{\oplus c} \to 0
$$
where $c = \dim((A/\pi A)_\mathfrak p)$.
\end{lemma}

\begin{proof}
By Lemma \ref{lemma-neron-functorial} we may replace $A$ by a localization
at an element not in $\mathfrak p$; we will use this without further mention.
Write $\kappa = R/\pi R$. Since smoothness is stable under base change
(Algebra, Lemma \ref{algebra-lemma-base-change-smooth})
we see that $A/\pi A$ is smooth over $\kappa$ at $\mathfrak p$.
Hence $(A/\pi A)_\mathfrak p$ is a regular local ring
(Algebra, Lemma \ref{algebra-lemma-characterize-smooth-over-field}).
Choose $g_1, \ldots, g_c \in \mathfrak p$ which map to
a regular system of parameters in $(A/\pi A)_\mathfrak p$.
Then we see that $\mathfrak p = (\pi, g_1, \ldots, g_c)$
after possibly replacing $A$ by a localization.
Note that $\pi, g_1, \ldots, g_c$ is a regular sequence
in $A_\mathfrak p$ (first $\pi$ is a nonzerodivisor and
then Algebra, Lemma \ref{algebra-lemma-regular-ring-CM}
for the rest of the sequence).
After replacing $A$ by a localization we may assume that
$\pi, g_1, \ldots, g_c$ is a regular sequence in $A$
(Algebra, Lemma \ref{algebra-lemma-regular-sequence-in-neighbourhood}).
It follows that
$$
A' = A[y_1, \ldots, y_c]/(\pi y_1 - g_1, \ldots, \pi y_c - g_c) =
A[y_1, \ldots, y_c]/I
$$
by More on Algebra, Lemma \ref{more-algebra-lemma-blowup-regular-sequence}.
In the following we will use the definition of smoothness using the
naive cotangent complex (Algebra, Definition \ref{algebra-definition-smooth})
and the criterion of Algebra, Lemma \ref{algebra-lemma-smooth-at-point}
without further mention.
The exact sequence of Algebra, Lemma \ref{algebra-lemma-exact-sequence-NL}
for $R \to A[y_1, \ldots, y_c] \to A'$ looks like this
$$
0 \to H_1(\NL_{A'/R}) \to I/I^2 \to
\Omega_{A/R} \otimes_A A' \oplus
\bigoplus\nolimits_{i = 1, \ldots, c} A' \text{d}y_i \to
\Omega_{A'/R} \to 0
$$
where the class of $\pi y_i - g_i$ in $I/I^2$ is mapped
to $- \text{d}g_i + \pi \text{d}y_i$ in the next term.
Here we have used Algebra, Lemma \ref{algebra-lemma-NL-surjection}
to compute $\NL_{A'/A[y_1, \ldots, y_c]}$ and we have used that
$R \to A[y_1, \ldots, y_c]$ is smooth, so
$H_1(\NL_{A[y_1, \ldots, y_c]/R}) = 0$ and
$\Omega_{A[y_1, \ldots, y_c]/R}$ is a finite projective (a fortiori flat)
$A[y_1, \ldots, y_c]$-module which is in fact the direct sum
of $\Omega_{A/R} \otimes_A A[y_1, \ldots, y_c]$ and a free
module with basis $\text{d}y_i$. To finish the proof it
suffices to show that $\text{d}g_1, \ldots, \text{d}g_c$
forms part of a basis for the finite free module $\Omega_{A/R, \mathfrak p}$.
Namely, this will show $(I/I^2)_\mathfrak p$ is free on $\pi y_i - g_i$,
the localization at $\mathfrak p$ of the middle map in the sequence is
injective, so $H_1(\NL_{A'/R})_\mathfrak p = 0$, and
that the cokernel $\Omega_{A'/R, \mathfrak p}$ is finite free.
To do this it suffices to show
that the images of $\text{d}g_i$ are $\kappa(\mathfrak p)$-linearly
independent in
$\Omega_{A/R, \mathfrak p}/\pi = \Omega_{(A/\pi A)/\kappa, \mathfrak p}$
(equality by Algebra, Lemma \ref{algebra-lemma-differentials-base-change}).
Since $\kappa \subset \kappa(\mathfrak p) \subset \Lambda/\pi \Lambda$
we see that $\kappa(\mathfrak p)$ is separable over $\kappa$
(Algebra, Definition \ref{algebra-definition-separable-field-extension}).
The desired linear independence now follows from
Algebra, Lemma \ref{algebra-lemma-computation-differential}.
\end{proof}

\end{document}
